
\section{定义的若干变体}\label{sec:7-2}
\zhindex{dingyideruoganbianti}{定义的若干变体}

本节首先讨论：当流形 $M$ 有一个自然的（例如可积的）、但非正则的\BookRecordIndex{T0088}殆复结构时，如何
计算\BookRecordIndex{T0143}格罗莫夫--威滕不变量。随后，我们为以固定树为模型或带\BookRecordIndex{T0150}固定标记点的\BookRecordIndex{T0340}稳定映射定义
一些相关不变量。

实际计算 $\operatorname{GW}_{A,k}$ 时，并不需要假设\BookRecordIndex{T0088}殆复结构 $J_0$ 正则，即属于
$\mathcal J_{\mathrm{reg}}(M,\omega)$。关键在于，\BookRecordIndex{T0283}求值映射
\[
  \operatorname{ev}_{J_0}:\mathcal M^*_{0,k}(A;J_0)\longrightarrow M^k
\]
应当是一个\BookRecordIndex{T0337}伪循环，并且表示与正则 $J$ 情形相同的同调类。为此，$J_0$ 只须满足下面
两个条件。

\begin{enumerate}[label=(\alph*),leftmargin=3em]
\item $J_0$ 对 $A$-曲线正则，即对每个表示类 $A$ 的 $J_0$-全纯球面，\BookRecordIndex{T0346}线性化算子
$D_u$ 都满射。于是对任意 $k$，\BookRecordIndex{T0246}模空间 $\mathcal M^*_{0,k}(A;J_0)$ 都有正确维数
\[
  d:=2n+2c_1(A)+2k-6.
\]
\item “边界”
\[
  \overline{\mathcal M}_{0,k}(A;J_0)
  \setminus\mathcal M^*_{0,k}(A;J_0)
\]
在 $\operatorname{ev}_{J_0}$ 下的像，其维数（在第~\ref{sec:6-5} 节的意义下）至多为
$d-2$。
\sourcefootnote{原书在条件 (b) 中把主\BookRecordIndex{T0072}层误写成
$\mathcal M^*_{0,k}(A;J)$，并在下文证明中重复了这一写法；但此处讨论的是
$\operatorname{ev}_{J_0}$ 的边界，且 $J$ 尚未指定，故两处均应为
$\mathcal M^*_{0,k}(A;J_0)$。}
\end{enumerate}
这些条件蕴含 $\operatorname{ev}_{J_0}$ 是一个 $d$ 维\BookRecordIndex{T0337}伪循环（见
定义~\ref{def:6-5-1}）。

\begin{lemma}\label{lem:7-2-1}
设 $J_0$ 是闭\BookRecordIndex{T0058}半正\BookRecordIndex{T0358}辛流形 $(M,\omega)$ 上被 $\omega$ \BookRecordIndex{T0372}驯服的\BookRecordIndex{T0088}殆复结构，并满足上述条件
(a)、(b)。则\BookRecordIndex{T0283}求值映射
\[
  \operatorname{ev}_{J_0}:\mathcal M^*_{0,k}(A;J_0)\longrightarrow M^k
\]
是由 $J\in\mathcal J_{\mathrm{reg}}(M,\omega)$ 时的
$\operatorname{ev}_J$ 所给\BookRecordIndex{T0145}格罗莫夫--威滕伪循环的\BookRecordIndex{T0292}弱代表元。
\end{lemma}

\begin{proof}
由引理~\ref{lem:6-5-7}，只须证明：对每个定义在维数为 $2nk-d$ 的紧定向光滑流形 $X$
上的光滑映射 $f:X\to M^k$，$\operatorname{ev}_{J_0}$ 与 $f$ 的交数都等于
$J\in\mathcal J_{\mathrm{reg}}(M,\omega)$ 时 $\operatorname{ev}_J$ 与 $f$ 的交数。
依假设，可以选取 $f:X\to M^k$，使得
\[
  \operatorname{ev}_{J_0}\Bigl(
  \overline{\mathcal M}_{0,k}(A;J_0)
  \setminus\mathcal M^*_{0,k}(A;J_0)\Bigr)
  \cap f(X)=\varnothing.
\]
设 $J^\nu\in\mathcal J_{\mathrm{reg}}(M,\omega)$ 收敛到 $J_0$。把\BookRecordIndex{T0380}隐函数定理应用于
$\operatorname{ev}_{J_0}$ 可知，对每个充分大的 $\nu$，
$\operatorname{ev}_{J^\nu}$ 与 $f$ 的交点数至少和
$\operatorname{ev}_{J_0}$ 与 $f$ 的交点数一样多。假设对所有充分大的 $\nu$，前者还有
更多交点；那么这些多出的交点必收敛到一个与 $\operatorname{im}f$ 相交的
$J_0$-全纯\BookRecordIndex{T0340}稳定映射，但由\BookRecordIndex{T0380}隐函数定理的唯一性部分，它不属于
$\operatorname{ev}_{J_0}(\mathcal M^*_{0,k}(A;J_0))$。这与 $f$ 避开
$\overline{\mathcal M}_{0,k}(A;J_0)\setminus\mathcal M^*_{0,k}(A;J_0)$ 的像矛盾。
因此，\BookRecordIndex{T0380}隐函数定理给出 $\operatorname{ev}_{J^\nu}$ 与 $f$ 的交点和
$\operatorname{ev}_{J_0}$ 与 $f$ 的交点之间的保定向一一对应。故充分大的 $\nu$ 满足
\[
  \operatorname{ev}_{J^\nu}\cdot f=\operatorname{ev}_{J_0}\cdot f,
\]
引理得证。
\end{proof}

\begin{example}\label{ex:7-2-2}
令 $M=S^2\times S^2$，记
\[
  A_1=[S^2\times\{\mathrm{pt}\}],
  \qquad
  A_2=[\{\mathrm{pt}\}\times S^2],
\]
取 $A$ 为对角类 $A_1+A_2$，并选择 $J_0$，使反对角类 $A_1-A_2$ 有一个嵌入
$J_0$-全纯代表元 $C$。由引理~\ref{lem:3-3-3}，$J_0$ 对类 $A$ 正则，故
$\mathcal M^*_{0,3}(A;J_0)$ 的维数为 $12$。反对角类没有正则代表元
（由引理~\ref{lem:3-3-2}），所以\BookRecordIndex{T0088}殆复结构 $J_0$ 不属于
$\mathcal J_{\mathrm{reg}}(M,\omega)$。另一方面，\BookRecordIndex{T0186}交点正性蕴含
$\mathcal M_{0,0}(A_1;J_0)$ 为空，并且 $C$ 唯一。因此，
$\overline{\mathcal M}_{0,3}(A;J_0)\setminus\mathcal M^*_{0,3}(A;J_0)$ 中仅有如下并集：
反对角曲线 $C$ 与类 $2A_2$ 的某个代表元之并；后者或者由两个不同球面组成，或者是一个
球面的二重覆盖。因此，它在 $M^3$ 中的像维数至多为 $10$。所以 $J_0$ 满足上述条件
(a)、(b)，可以用它计算 $\operatorname{GW}_{A,3}$。
\end{example}

\paragraph{\BookRecordIndex{T0399}障碍丛。}
即使 $J_0$ 不满足定理~\ref{thm:7-1-1} 的\BookRecordIndex{T0408}正则性条件，有时仍可通过研究非正则\BookRecordIndex{T0246}模空间
$\mathcal M(A;J_0)$ 来计算 $A$ 的\BookRecordIndex{T0143}格罗莫夫--威滕不变量。微分拓扑中有一个类似情形：
利用向量丛的函子性质来计算其欧拉类，而不是计数横截截面的零点。我们只讨论一种非常
简单的情形：不带参数的 $A$-曲线\BookRecordIndex{T0246}模空间
\[
  \widehat{\mathcal M}_{J_0}:=\mathcal M(A;J_0)/G
\]
是光滑紧流形，并且 $\mathcal M(A;J_0)$ 在每条曲线 $u$ 处的切空间等于 $D_u$ 的核。
这时，各余核秩恒定，并拼成 $\widehat{\mathcal M}_{J_0}$ 上的一个向量丛，称为\BookRecordIndex{T0399}障碍丛。
如果这个向量丛的秩还等于 $\widehat{\mathcal M}_{J_0}$ 的维数，那么问题的指标为零，
也就是说，人们预期只有有限条 $A$-曲线。下面的结果说明，所得不变量恰为\BookRecordIndex{T0399}障碍丛的欧拉数。
更精确地说，我们考虑如下情形。

\medskip
\noindent\textbf{(H)}
\emph{设 $(M,\omega)$ 是紧\BookRecordIndex{T0058}半正辛 $2n$ 维流形，$J_0$ 是 $M$ 上被 $\omega$ \BookRecordIndex{T0372}驯服的\BookRecordIndex{T0088}殆复
结构，而 $A\in H_2(M;\mathbb Z)$ 是满足
\[
  2n+2c_1(A)-6=0
\]
的同调类。假设每个表示类 $A$ 的 $J_0$-全纯\BookRecordIndex{T0340}稳定映射都是单纯 $J_0$-全纯球面，
$\mathcal M(A;J_0)$ 是 $C^\infty(S^2,M)$ 中可定向的光滑子流形，商空间
$\widehat{\mathcal M}_{J_0}:=\mathcal M(A;J_0)/G$ 紧，并且对每个
$u\in\mathcal M(A;J_0)$ 都有
\[
  T_u\mathcal M(A;J_0)=\ker D_u.
\]}

在这些假设下，群 $G=\operatorname{PSL}(2,\mathbb C)$ 在 $\mathcal M(A;J_0)$ 上自由
作用，所以 $\widehat{\mathcal M}_{J_0}$ 是光滑紧流形。令
$E_{J_0}\to\mathcal M(A;J_0)$ 表示可定向实向量丛，其纤维为
\begin{equation}\label{eq:7-2-1}
  E_{u,J_0}:=(\operatorname{im}D_u)^\perp
  \cong\operatorname{coker}D_u
  =\Omega^{0,1}(S^2,u^*TM)/\operatorname{im}D_u.
\end{equation}
这里的正交补相对于 $\Omega^{0,1}(S^2,u^*TM)$ 上的 $L^2$ 内积；该内积由 $M$ 上与
$\omega$、$J_0$ 相联系的黎曼度量决定。由于一形式空间上的内积只依赖于 $S^2$ 上的
复结构，而不依赖于体积形式，向量丛 $E_{J_0}$ 在
$G=\operatorname{PSL}(2,\mathbb C)$ 的作用下不变（见注~\ref{rem:3-1-2}）。由于 $D_u$ 的
\BookRecordIndex{T0129}弗雷德霍姆指标为 $2n+2c_1(A)=6$，商丛
\[
  \widehat E_{J_0}:=E_{J_0}/G\longrightarrow\widehat{\mathcal M}_{J_0}
\]
的秩为
\[
  \operatorname{rank}\widehat E_{J_0}
  =\dim\operatorname{coker}D_u
  =\dim\ker D_u-6
  =\dim\widehat{\mathcal M}_{J_0}=:d.
\]
向量丛 $\widehat E_{J_0}\to\widehat{\mathcal M}_{J_0}$ 称为\emph{\BookRecordIndex{T0399}障碍丛}。

\begin{theorem}[\BookRecordIndex{T0399}障碍丛]\label{thm:7-2-3}
假设 (H) 成立。则\BookRecordIndex{T0143}格罗莫夫--威滕不变量等于\BookRecordIndex{T0399}障碍丛的欧拉数：
\begin{equation}\label{eq:7-2-2}
  \operatorname{GW}^{M}_{A,0}
  =\int_{\widehat{\mathcal M}_{J_0}}e(\widehat E_{J_0}).
\end{equation}
\end{theorem}

为理解这个公式，注意向量丛 $E_{J_0}\to\mathcal M(A;J_0)$ 的总空间在零截面元素处的
切空间，是底空间切空间与 $E_{J_0}$ 纤维的直和；也就是说，对每个
$u\in\mathcal M(A;J_0)$，
\[
  T_{(u,0)}E_{J_0}=\ker D_u\oplus\operatorname{coker}D_u.
\]
因此，$D_u$ 的\BookRecordIndex{T0162}行列式线的定向确定了\BookRecordIndex{T0399}障碍丛总空间 $\widehat E_{J_0}$ 的定向。虽然\BookRecordIndex{T0246}模空间
$\mathcal M(A;J_0)$ 未必有自然定向，我们可以任意固定它的一个定向，从而也固定
$\widehat{\mathcal M}_{J_0}$ 的定向。反转 $\widehat{\mathcal M}_{J_0}$ 的定向时，
$\widehat E_{J_0}$ 各纤维的定向也随之反转。因此，式~\eqref{eq:7-2-2} 的右端保持不变。

恒等式~\eqref{eq:7-2-2} 与 Cieliebak--Mundet--Salamon \srcite{66} 关于具有紧零点集的
Hilbert 空间丛之弗雷德霍姆截面的（等变）欧拉类的结果密切相关。不过，他们的论证不能
直接应用于当前情形，因为他们假设\BookRecordIndex{T0413}重参数化群 $G$ 紧，并且在相关 Hilbert 空间丛上光滑
作用。为绕开这个问题，我们采用\BookRecordIndex{T0265}平衡映射 $u:S^2\to M$。

\begin{exercise}\label{ex:7-2-4}
令 $x_1,x_2,x_3:S^2\to\mathbb R$ 表示 $S^2\subset\mathbb R^3$ 上的标准坐标函数。如果
二形式 $\sigma\in\Omega^2(S^2)$ 满足
\begin{equation}\label{eq:7-2-3}
  \int_{S^2}x_i\sigma=0,
  \qquad i=1,2,3,
\end{equation}
就称它是\emph{平衡的}。这等价于说，把 $(S^2,\sigma)$ 看成 $\mathbb R^3$ 的子集时，
其质心位于原点。若光滑映射 $u:S^2\to M$ 使 $u^*\omega$ 是 $S^2$ 上的平衡二形式，
就称 $u$ 是\emph{平衡的}。设 $\sigma\in\Omega^2(S^2)$ 且
$\int_{S^2}\sigma\ne0$。证明：存在 $\phi\in\operatorname{PSL}(2,\mathbb C)$，使
$\phi^*\sigma$ 平衡；在右侧乘以 $\operatorname{SO}(3)$ 的元素的意义下，这样的 $\phi$
唯一。（见 Cieliebak--Gaio--Salamon \srcite{64}，命题 3.3。）
\end{exercise}

\begin{proof}[定理~\ref{thm:7-2-3} 的证明]
证明分六步进行。

\medskip
\noindent\textbf{步骤 1（平衡条件）。}\quad
令 $\mathcal B\subset C^\infty(S^2,M)$ 为表示同调类 $A$ 的光滑映射
$u:S^2\to M$ 所成的空间，并令
\[
  \mathcal B_0:=\bigl\{u\in C^\infty(S^2,M)\mid
  [u]=A\text{ 且 }u^*\omega\text{ 满足式~\eqref{eq:7-2-3}}\bigr\}
\]
为表示类 $A$ 的\BookRecordIndex{T0265}平衡映射之空间。则 $\mathcal B_0$ 是 $\mathcal B$ 中光滑的余维三
弗雷歇子流形，其在 $u$ 处的切空间由所有沿 $u$ 的向量场
$\xi\in\Omega^0(S^2,u^*TM)$ 组成，并满足
\begin{equation}\label{eq:7-2-4}
  \int_{S^2}dx_i\wedge\omega\bigl(\xi,du(\,\cdot\,)\bigr)=0,
  \qquad i=1,2,3.
\end{equation}
包含映射 $\mathcal B_0\hookrightarrow\mathcal B$ 诱导双射
\[
  \mathcal B_0/\operatorname{SO}(3)
  \cong\mathcal B/\operatorname{PSL}(2,\mathbb C).
\]

\BookRecordIndex{T0130}弗雷歇流形 $\mathcal B$ 带有\BookRecordIndex{T0363}辛形式
\[
  (\xi,\eta)\longmapsto
  \int_{S^2}\omega(\xi,\eta)\,\mathrm{dvol}_{S^2},
  \qquad
  \xi,\eta\in\Omega^0(S^2,u^*TM)=T_u\mathcal B.
\]
$\operatorname{SO}(3)$ 在 $\mathcal B$ 上的作用是哈密顿群作用，由等变\BookRecordIndex{T0204}矩映射
\[
  \mathcal B\longrightarrow\mathbb R^3:
  u\longmapsto
  \left(\int_{S^2}x_1u^*\omega,
        \int_{S^2}x_2u^*\omega,
        \int_{S^2}x_3u^*\omega\right)
\]
生成。由于 $\omega(A)>0$，习题~\ref{ex:7-2-4} 蕴含：对每个 $u\in\mathcal B$，都存在
\BookRecordIndex{T0040}Möbius 变换 $\phi\in G$，使 $u\circ\phi$ 平衡；在右侧复合一个保定向等距同构的意义下，
$\phi$ 唯一。此外，由于 $\operatorname{SO}(3)$ 在 $\mathcal B$ 上的作用只有有限\BookRecordIndex{T0342}稳定
子群，零是这个\BookRecordIndex{T0204}矩映射的\BookRecordIndex{T0409}正则值，故 $\mathcal B_0$ 是 $\mathcal B$ 的余维三子流形。这就
证明了步骤 1。

\medskip
\noindent\textbf{步骤 2。}\quad
表示类 $A$ 的平衡 $J_0$-全纯曲线之空间
\[
  \mathcal M_{J_0}:=\mathcal M(A;J_0)\cap\mathcal B_0
\]
是 $\mathcal B_0$ 中维数为 $d+3$ 的光滑紧子流形。群 $\operatorname{SO}(3)$ 在
$\mathcal M_{J_0}$ 上自由作用，而且
\[
  \widehat{\mathcal M}_{J_0}
  =\mathcal M_{J_0}/\operatorname{SO}(3).
\]
$E_{J_0}$ 在 $\mathcal M_{J_0}$ 上的限制（仍记为 $E_{J_0}$）是
$\operatorname{SO}(3)$-等变向量丛，而\BookRecordIndex{T0399}障碍丛
$\widehat E_{J_0}\to\widehat{\mathcal M}_{J_0}$ 是商丛
$E_{J_0}/\operatorname{SO}(3)$。

步骤 2 立即来自假设 (H) 与步骤 1。

\medskip
\noindent\textbf{步骤 3（延拓\BookRecordIndex{T0399}障碍丛）。}\quad
固定常数 $p>2$，简记 $\mathcal J:=\mathcal J_\tau(M,\omega)$，并令
$\mathcal E\to\mathcal B\times\mathcal J$ 表示纤维为
\[
  \mathcal E_{u,J}:=\Omega_J^{0,1}(S^2,u^*TM)
\]
的向量丛。则存在 $\mathcal M_{J_0}$ 在 $\mathcal B_0$ 中的一个
$\operatorname{SO}(3)$-不变邻域 $\mathcal U_0\subset\mathcal B_0$，它在
$\mathcal B_0$ 的 $W^{1,p}$ 拓扑中开；还存在 $J_0$ 在 $\mathcal J$ 中的一个邻域
$\mathcal V_0\subset\mathcal J$，它在 $\mathcal J$ 的 $C^0$ 拓扑中开；使得向量丛
$E_{J_0}\to\mathcal M_{J_0}$ 延拓为 $\mathcal E$ 在
$\mathcal U_0\times\mathcal V_0$ 上的一个 $\operatorname{SO}(3)$-等变实秩 $d$ 子丛。
把这个延拓子丛记作
\[
  E\longrightarrow\mathcal U_0\times\mathcal V_0.
\]

对 $u\in\mathcal B$ 和 $J\in\mathcal J$，令
\[
  \pi_{u,J}:\Omega^1(S^2,u^*TM)
  \longrightarrow\Omega_J^{0,1}(S^2,u^*TM)
\]
表示自然投影。于是映射 $(u,\eta)\mapsto\pi_{u,J}\eta$ 是 $G$-等变的。给 $M$ 配备由
$\omega$、$J_0$ 确定的黎曼度量，并用 $\mathcal N$ 表示
$\mathcal M_{J_0}$ 在 $\mathcal B$ 中的\BookRecordIndex{T0117}法丛。$\mathcal N$ 的元素是偶对 $(u,\xi)$，
其中 $u\in\mathcal M_{J_0}$，而
$\xi\in\Omega^0(S^2,u^*TM)=T_u\mathcal B$ 在 $L^2$ 意义下正交于
$D_u$ 的核与 $T_u\mathcal B_0$ 的交；即
\[
  D_u\xi_0=0,\quad
  \xi_0\text{ 满足式~\eqref{eq:7-2-4}}
  \quad\Longrightarrow\quad
  \int_{S^2}\langle\xi,\xi_0\rangle\,\mathrm{dvol}_{S^2}=0.
\]
\BookRecordIndex{T0117}法丛 $\mathcal N$ 在 $\operatorname{SO}(3)$ 的作用下不变，而且映射
\begin{equation}\label{eq:7-2-5}
  \mathcal N\longrightarrow\mathcal B:
  (u,\xi)\longmapsto\exp_u(\xi)
\end{equation}
是 $\operatorname{SO}(3)$-等变的。需要强调，以协变导数定义的 $\xi$ 的
$W^{1,p}$ 范数，在 $\mathcal N$ 上的 $\operatorname{SO}(3)$ 作用下不变。\BookRecordIndex{T0380}隐函数定理
断言，当 $\delta>0$ 充分小时，指数映射~\eqref{eq:7-2-5} 在集合
\[
  \mathcal N_\delta:=
  \bigl\{(u,\xi)\in\mathcal N\mid
  \|\xi\|_{L^p}+\|\nabla\xi\|_{L^p}<\delta\bigr\}
\]
上的限制，是从 $\mathcal N_\delta$ 到 $\mathcal M_{J_0}$ 的某个
$\operatorname{SO}(3)$-不变邻域 $\mathcal U_\delta\subset\mathcal B$ 的微分同胚；
$\mathcal U_\delta$ 在 $W^{1,p}$ 拓扑下开。对
$(u,\xi)\in\mathcal N_\delta$，用
\[
  \Phi(u,\xi):\Omega^1(S^2,u^*TM)
  \longrightarrow\Omega^1(S^2,\exp_u(\xi)^*TM)
\]
表示沿测地线 $t\mapsto\exp_{u(z)}(t\xi(z))$ \BookRecordIndex{T0263}平行移动所给的同构。若 $\delta>0$ 充分小，
且 $J\in\mathcal J$ 在 $C^0$ 拓扑中充分接近 $J_0$，那么对每个
$(u,\xi)\in\mathcal N_\delta$，线性映射
\[
  \pi_{\exp_u(\xi),J}\Phi(u,\xi):
  \Omega^1(S^2,u^*TM)
  \longrightarrow
  \Omega_J^{0,1}(S^2,\exp_u(\xi)^*TM)
\]
在子空间 $E_{u,J_0}$ 上的限制是单射
\sourcefootnote{原书把这个线性映射两端的一形式空间底空间都印作 $M$；但
$u:S^2\to M$，且紧邻定义的 $\Phi(u,\xi)$ 与 $\pi_{u,J}$ 都作用在 $S^2$ 上的形式，
故两端唯一应以 $S^2$ 为底。}。因此，存在 $J_0$ 在 $\mathcal J$ 中的一个 $C^0$ 开邻域
$\mathcal V_0$ 以及 $\delta>0$，使公式
\begin{equation}\label{eq:7-2-6}
  E_{\exp_u(\xi),J}
  :=\pi_{\exp_u(\xi),J}\Phi(u,\xi)E_{u,J_0},
  \qquad
  (u,\xi)\in\mathcal N_\delta,
  \quad J\in\mathcal V_0,
\end{equation}
定义 $\mathcal U_\delta\times\mathcal V_0$ 上一个
$\operatorname{SO}(3)$-等变光滑秩 $d$ 向量丛 $E$，它在
$\mathcal M_{J_0}\times\{J_0\}$ 上的限制与 $E_{J_0}$ 一致。定义
$\mathcal U_0:=\mathcal U_\delta\cap\mathcal B_0$，并把 $E$ 限制到乘积
$\mathcal U_0\times\mathcal V_0$。这个向量丛满足步骤 3 的要求。

\medskip
\noindent\textbf{步骤 4（关键估计）。}\quad
固定常数 $p>2$，令 $E\to\mathcal U_0\times\mathcal V_0$ 为步骤 3 中构造的
$\mathcal E$ 的子丛~\eqref{eq:7-2-6}。则对每个整数 $k\geq0$ 和每个
$(u,J)\in\mathcal U_0\times\mathcal V_0$，都存在常数
$c_k=c_k(u,J)>0$，使得
\begin{equation}\label{eq:7-2-7}
  \|\eta\|_{W^{k+1,p}}
  \leq c_k(u,J)\|\eta\|_{W^{k,p}}
  \qquad\text{对所有 }\eta\in E_{u,J}.
\end{equation}
常数 $c_k$ 只依赖于 $u$ 的 $W^{k+1,p}$ 范数和 $J$ 的 $C^{k+1}$ 范数。当
$J=J_0$ 且 $u\in\mathcal M_{J_0}$ 时，这一点显然成立：对每个这样的 $u$，
$E_{u,J_0}$ 都是由光滑 $(0,1)$-形式组成的有限维向量空间，而
$\mathcal M_{J_0}$ 紧。一般情形下，取
$(\exp_u(\xi),J)\in\mathcal U_0\times\mathcal V_0$，考虑同构
\[
  \pi_{\exp_u(\xi),J}\Phi(u,\xi):
  E_{u,J_0}\longrightarrow E_{\exp_u(\xi),J}.
\]
这个同构及其逆在两端 $W^{k+1,p}$ 范数下的算子范数，由 $\xi$ 的
$W^{k+1,p}$ 范数和 $J$ 的 $C^{k+1}$ 范数控制。这就证明了步骤 4。

\medskip
\noindent\textbf{步骤 5（选择邻域）。}\quad
令 $E\to\mathcal U_0\times\mathcal V_0$ 为步骤 3 中构造的
$\mathcal E$ 的秩 $d$ 子丛。存在 $\mathcal M_{J_0}$ 在 $\mathcal U_0$ 中的邻域
$\mathcal U\subset\mathcal U_0$，以及 $J_0$ 在 $\mathcal V_0$ 中的邻域
$\mathcal V\subset\mathcal V_0$，二者均在 $C^\infty$ 拓扑下开，并满足下列性质。
分别用 $\overline{\mathcal U}$ 和 $\overline{\mathcal V}$ 表示
$\mathcal U$、$\mathcal V$ 的 $C^\infty$ 闭包。
\begin{enumerate}[label=(\alph*),leftmargin=3em]
\item 向量丛 $E$ 延拓到
$\overline{\mathcal U}\times\overline{\mathcal V}$ 上，$\mathcal V$ 连通，且
$\operatorname{SO}(3)$ 在 $\overline{\mathcal U}$ 上自由作用。
\item 对 $J\in\overline{\mathcal V}$ 定义
\begin{equation}\label{eq:7-2-8}
  \mathcal M_{J,E}:=
  \bigl\{u\in\overline{\mathcal U}\mid
  \bar\partial_J(u)\in E_{u,J}\bigr\}\subset\mathcal B_0.
\end{equation}
则 $\mathcal M_{J_0,E}=\mathcal M_{J_0}$，并且对每个
$J\in\overline{\mathcal V}$ 都有 $\mathcal M_{J,E}\subset\mathcal U$。
\item 对每个紧集 $K\subset\overline{\mathcal V}$，集合
$\bigcup_{J\in K}\mathcal M_{J,E}\subset\mathcal U$ 都是紧的。
\item 对每个 $J\in\overline{\mathcal V}$，令
$\mathcal E_J\to\overline{\mathcal U}$ 为纤维
\[
  \mathcal E_{u,J}:=\Omega_J^{0,1}(S^2,u^*TM)
\]
的无限维向量丛，并令 $E_J\to\overline{\mathcal U}$ 为纤维 $E_{u,J}$ 的秩 $d$
子丛。则截面
\[
  \overline{\mathcal U}\longrightarrow\mathcal E_J,
  \qquad u\longmapsto\bar\partial_J(u),
\]
与 $E_J$ 横截。
\item 若 $J\in\overline{\mathcal V}$，则每个表示类 $A$ 的 $J$-全纯\BookRecordIndex{T0340}稳定映射，都是
落在 $\mathcal U$ 中的单纯 $J$-全纯球面。
\end{enumerate}

断言 (a) 是显然的。先证明可以选择 $\mathcal U$，使 (b)、(d) 对 $J=J_0$ 成立。
固定满足 (a) 的开集 $\mathcal U\subset\mathcal U_0$ 与
$\mathcal V\subset\mathcal V_0$。由式~\eqref{eq:7-2-1}，沿子流形
$\mathcal M_{J_0}\subset\mathcal M_{J_0,E}$，子丛
$E_{J_0}\subset\mathcal E_{J_0}$ 与截面
\[
  \overline{\mathcal U}\longrightarrow\mathcal E_{J_0},
  \qquad u\longmapsto\bar\partial_{J_0}(u)
\]
横截。在适当的 Banach 流形框架中应用\BookRecordIndex{T0380}隐函数定理可知，$\mathcal M_{J_0}$ 是
$\mathcal M_{J_0,E}$ 的孤立子集；换言之，$\mathcal M_{J_0}$ 在
$\overline{\mathcal U}$ 中有一个 $W^{1,p}$ 开邻域，其中没有
$\mathcal M_{J_0,E}$ 的其他元素。必要时缩小 $\mathcal U$，便得到
$\mathcal M_{J_0,E}=\mathcal M_{J_0}$。所以 (b)、(d) 对 $J=J_0$ 成立。

接着证明可选择 $\mathcal V$，使 (b) 对每个 $J\in\overline{\mathcal V}$ 成立。
反设对 $J_0$ 的每个 $C^\infty$ 邻域，(b) 都不成立。则存在序列
$J^\nu\in\mathcal J$ 和 $u^\nu\in\overline{\mathcal U}\setminus\mathcal U$，使
$J^\nu$ 在 $C^\infty$ 拓扑下收敛到 $J_0$，且
$\bar\partial_{J^\nu}(u^\nu)\in E_{u^\nu,J^\nu}$。按假设，$u^\nu$ 在
$W^{1,p}$ 中一致有界。由步骤 4 的估计~\eqref{eq:7-2-7} 和\BookRecordIndex{T0330}椭圆自举，$u^\nu$ 对每个
$k\geq0$ 都在 $W^{k,p}$ 中一致有界。因此，$u^\nu$ 有一个在 $C^\infty$ 拓扑下
收敛的子列，其极限 $u^0\in\overline{\mathcal U}\setminus\mathcal U$ 满足
$\bar\partial_{J_0}(u^0)\in E_{u^0,J_0}$；这与
$\mathcal M_{J_0,E}=\mathcal M_{J_0}$ 矛盾。故可选择 $\mathcal V$，使 (b) 对每个
$J\in\overline{\mathcal V}$ 都成立。

现在证明 $\mathcal U$、$\mathcal V$ 满足 (c)。设
$K\subset\overline{\mathcal V}$ 紧，且 $J^\nu\in K$、
$u^\nu\in\mathcal M_{J^\nu,E}$。取子列后可设 $J^\nu$ 在 $C^\infty$ 拓扑下收敛于
$J\in K$。再次使用式~\eqref{eq:7-2-7} 和\BookRecordIndex{T0330}椭圆自举可知，$u^\nu$ 对每个整数
$k\geq0$ 都在 $W^{k,p}$ 中一致有界。进一步取子列后，不妨设 $u^\nu$ 在
$C^\infty$ 拓扑下收敛于某个 $u\in\overline{\mathcal U}$。极限满足
$\bar\partial_J(u)\in E_{u,J}$，故属于
$\mathcal M_{J,E}\subset\mathcal U$。这就证明了 (c)。

必要时进一步缩小 $\mathcal V$。由于\BookRecordIndex{T0166}横截性是开条件，可得到 (d)；由\BookRecordIndex{T0139}格罗莫夫紧性和
假设 (H)，可得到 (e)。步骤 5 得证。

\medskip
\noindent\textbf{步骤 6。}\quad
下面证明定理。令 $E\to\mathcal U_0\times\mathcal V_0$ 为步骤 3 中构造的
$\mathcal E$ 的秩 $d$ 子丛~\eqref{eq:7-2-6}，并取步骤 5 中的
$\mathcal U\subset\mathcal U_0$ 与 $\mathcal V\subset\mathcal V_0$。对
$J\in\mathcal V$，令 $\mathcal M_{J,E}\subset\mathcal B_0$ 由
式~\eqref{eq:7-2-8} 定义。由 (b)、(c)、(d)，对每个 $J\in\mathcal V$，
$\mathcal M_{J,E}$ 都是无边界的 $d+3$ 维紧流形。它还带有
$\operatorname{SO}(3)$ 的光滑自由作用。秩 $d$ 向量丛
$E_J\to\mathcal M_{J,E}$ 是 $\operatorname{SO}(3)$-等变的，因而下降为商丛
\[
  \widehat E_J:=E_J/\operatorname{SO}(3)
  \longrightarrow
  \widehat{\mathcal M}_{J,E}:=
  \mathcal M_{J,E}/\operatorname{SO}(3).
\]
由 (a)、(c)，$\widehat E_J\to\widehat{\mathcal M}_{J,E}$ 的欧拉数与 $J$ 无关。
此外，对每个 $J\in\mathcal V$，存在光滑的 $\operatorname{SO}(3)$-等变截面
\[
  S_J:\mathcal M_{J,E}\longrightarrow E_J,
  \qquad S_J(u):=\bar\partial_J(u)\in E_{u,J},
\]
它在 $J=J_0$ 时恒为零。现在取 $J\in\mathcal V$ 为正则\BookRecordIndex{T0088}殆复结构，即对每个
$u\in\mathcal M(A;J)$，$D_{u,J}$ 都满射（这样的 $J$ 由
定理~\ref{thm:3-1-6} 存在）。则诱导截面
\[
  \widehat S_J:\widehat{\mathcal M}_{J,E}\longrightarrow\widehat E_J
\]
与零截面横截；由 (e)，其带符号计数的零点数正是所求\BookRecordIndex{T0143}格罗莫夫--威滕不变量。因此
\[
  \int_{\widehat{\mathcal M}_{J_0}}e(\widehat E_{J_0})
  =\int_{\widehat{\mathcal M}_{J_0,E}}e(\widehat E_{J_0})
  =\int_{\widehat{\mathcal M}_{J,E}}e(\widehat E_J)
  =\operatorname{GW}^{M}_{A,0}.
\]
这就证明了定理~\ref{thm:7-2-3}。
\end{proof}

下面的例子说明如何使用公式~\eqref{eq:7-2-2}。

\begin{example}\label{ex:7-2-5}
设 $L\to S^2\times S^2$ 是双次数为 $(1,-2)$ 的线丛。换言之，令
\[
  B_1:=[S^2\times\{\mathrm{pt}\}],
  \qquad
  B_2:=[\{\mathrm{pt}\}\times S^2],
\]
则
\[
  \langle c_1(L),B_1\rangle=1,
  \qquad
  \langle c_1(L),B_2\rangle=-2.
\]
把 $(M,J_0)$ 定义为射影化
\[
  (M,J_0):=(\mathbb P(L\oplus\mathbb C),J_0).
\]
具体地说，$M$ 是由 $T^3=S^1\times S^1\times S^1$ 在 $\mathbb C^6$ 上的如下作用
得到的约化空间：
\[
  (\lambda_1,\lambda_2,\lambda_3)\cdot
  (z_0,z_1,w_0,w_1,v_0,v_1)
  :=(\lambda_1z_0,\lambda_1z_1,
     \lambda_2w_0,\lambda_2w_1,
     \lambda_1\lambda_2^{-2}\lambda_3v_0,\lambda_3v_1).
\]
这个作用的\BookRecordIndex{T0204}矩映射为
\[
  \mu:(z,w,v)\longmapsto
  \bigl(|z|^2+|v_0|^2,\ |w|^2-2|v_0|^2,\ |v|^2\bigr),
\]
其中 $|z|^2:=|z_0|^2+|z_1|^2$，其余类似。定义
\[
  M:=\mu^{-1}(3,1,1)/T^3.
\]
于是 $M$ 是环面\BookRecordIndex{T0358}辛流形。注意其\BookRecordIndex{T0363}辛形式 $\omega$ 是整的，并且 $H_2(M;\mathbb Z)$ 由
下列球面生成：
\[
\begin{aligned}
  [z_0:z_1]&\longmapsto[z_0:z_1:0:1:0:1],\\
  [w_0:w_1]&\longmapsto[3:0:w_0:w_1:0:1],\\
  [v_0:v_1]&\longmapsto
  [0:\sqrt{3-|v_0|^2}:0:\sqrt{1+2|v_0|^2}:v_0:v_1].
\end{aligned}
\]
它们的辛面积依次为 $3,1,1$。

$S^2\times S^2$ 中的 $B_2$-球面 $w\mapsto(z,w)$ 提升为 $M$ 中由
\[
  u_z(w)=[z_0:z_1:w_0:w_1:0:1]
\]
给出的一族球面。这些球面都位于\BookRecordIndex{T0279}纤维化 $M\to S^2\times S^2$ 的截面
$V_-:=\{v_0=0\}$ 中，该截面的\BookRecordIndex{T0117}法丛是 $L$。因此，若 $A$ 表示这些提升球面所代表的类，
则
\[
  A\cdot[V_-]=-2,
  \qquad
  \omega(A)=1.
\]
截面 $V_-$ 是全纯的，故\BookRecordIndex{T0186}交点正性（见习题~\ref{ex:2-6-1}）蕴含：$M$ 中仅有的
$J_0$-全纯 $A$-球面，正是上述完全落在 $V_-$ 中的球面。因此，这些球面的\BookRecordIndex{T0246}模空间
\[
  \widehat{\mathcal M}:=\mathcal M(A;J_0)/G
\]
紧，并可等同于 $S^2$。另一方面，由于 $L$ 在每个 $A$-球面上的限制陈数为 $-2$，
$c_1(A)=0$；我们处于例~\ref{ex:7-1-4} 中 $\mu(A,0)=0$ 的情形。因此目标是计算类
$A$ 的 $J$-全纯代表元之数 $\operatorname{GW}^{M}_{A,0}$。

$M$ 上的自然复结构 $J_0$ 不是正则的。不过，对每个 $u\in\mathcal M(A;J_0)$，曲线
$C_z:=\operatorname{im}u_z$ 的\BookRecordIndex{T0117}法丛分解为
\[
  L_0\oplus L|_{C_z}\cong L_0\oplus L_{-2},
\]
这里 $L_k$ 表示 $S^2$ 上陈数为 $k$ 的全纯线丛。因此，$D_u$ 的核是
$L_2\oplus L_0$ 的全纯截面空间，其中把 $L_2$ 看作 $C_z$ 的切丛。于是
$\ker D_u$ 可与 $T_u\mathcal M(A;J_0)$ 等同，正如定理~\ref{thm:7-2-3} 所要求的；
我们可以利用\BookRecordIndex{T0399}障碍丛 $E$ 计算 $\operatorname{GW}^{M}_{A,0}$。

按定义，$E$ 在 $u$ 处的纤维就是 $\operatorname{coker}D_u$。由
引理~\ref{lem:3-3-1}，若 $u=u_z$，则可以把对偶空间
$(\operatorname{coker}D_u)^*$ 等同于全纯截面空间
\[
  H^0\bigl(S^2,L^*|_{C_z}\otimes K_{C_z}\bigr),
\]
其中 $L^*|_{C_z}\otimes K_{C_z}$ 是平凡线丛，而 $K_C=T^*C$ 表示 $C$ 的\BookRecordIndex{T0103}典范丛。
因此，对每个 $u$，向量空间 $(\operatorname{coker}D_u)^*$ 的复维数都是 $1$。各线丛
$K_{C_z}$ 拼成 $V_-=S^2\times S^2$ 上双次数为 $(0,-2)$ 的全纯线丛
$K\to V_-$，相应的空间 $(\operatorname{coker}D_u)^*$ 也拼成
$\widehat{\mathcal M}$ 上一个全纯线丛 $\widehat E^*\to\widehat{\mathcal M}$ 的纤维。
又因为 $L^*\otimes K\to V_-$ 同构于第一因子投影所拉回的
$L_{-1}\to S^2$，且 $(\operatorname{coker}D_u)^*$ 中的元素是常值截面，所以
$\widehat E^*$ 在 $u_z$ 处的纤维可与 $L_{-1}$ 在 $z$ 处的纤维
$L_z^{-1}$ 等同。因此，\BookRecordIndex{T0399}障碍丛
$\widehat E\to\widehat{\mathcal M}$ 可等同于 $S^2$ 上欧拉数（亦即陈数）为 $1$ 的线丛
$L_1\to S^2$。所以由式~\eqref{eq:7-2-2}，
\[
  \operatorname{GW}^{M}_{A,0}=1.
\]
\end{example}

\begin{exercise}\label{ex:7-2-6}
把定理~\ref{thm:7-2-3} 的结果推广到如下情形：
$2n+2c_1(A)-6=\mu>0$；\BookRecordIndex{T0246}模空间 $\mathcal M(A;J_0)$ 是维数为
$\mu+6+d>\mu+6$ 的光滑流形，其切丛由各核 $\ker D_u$ 组成；并且商空间
$\mathcal M(A;J_0)/G$ 紧。
\end{exercise}

上一习题的结果仍然十分受限。在许多有意义的情形中，虽然
$\mathcal M(A;J_0)/G$ 是\BookRecordIndex{T0074}超额维数 $d>0$ 的光滑流形，却并不紧。此时，理解\BookRecordIndex{T0399}障碍丛向紧化
的延拓结构非常重要，因为边界上的扭曲可能改变其欧拉类，因而也可能改变相应的
\BookRecordIndex{T0143}格罗莫夫--威滕不变量。代数几何学家已经发展出许多处理这一问题的精细方法；在此讨论
会使我们离题太远。有兴趣的读者可以先参阅 Okounkov--Pandharipande
\srcite{311,312}、Zinger \srcite{426,427,428} 及其中的参考文献。还要注意，代数几何中
所关心的枚举不变量并不总与\BookRecordIndex{T0143}格罗莫夫--威滕不变量相同；比较例~\ref{ex:8-6-12}。

本节最后讨论几种略有不同的不变量，它们计数具有特殊性质的曲线或\BookRecordIndex{T0340}稳定映射。

\paragraph{\BookRecordIndex{T0303}树不变量。}
设 $(M,\omega)$ 是紧\BookRecordIndex{T0058}半正\BookRecordIndex{T0358}辛流形。固定一棵 $k$-\BookRecordIndex{T0084}带标号树
$T=(T,E,\Lambda)$，以及一族球面同调类
$\{A_\gamma\}_{\gamma\in T}\subset H_2(M;\mathbb Z)$，它们满足第~\ref{sec:6-1} 节
的稳定性条件~\eqref{eq:6-1-1} 和定理~\ref{thm:6-6-1} 中的条件~\eqref{eq:6-6-1}。
（例如，若 $(M,\omega)$ 单调，则每个非零 $A_\gamma$ 都满足
式~\eqref{eq:6-6-1}。）由习题~\ref{ex:6-6-3}，对每个
$J\in\mathcal J_{\mathrm{reg}}(M,\omega)$，\BookRecordIndex{T0283}求值映射
\[
  \operatorname{ev}_{T,J,\{A_\gamma\}}:
  \mathcal M^*_{0,T}(\{A_\gamma\};J)\longrightarrow M^k
\]
都是维数为 $\mu(A,k)-2e(T)$ 的\BookRecordIndex{T0337}伪循环，其中
\[
  A:=\sum_\gamma[A_\gamma].
\]
因此可以定义不变量
\[
  \operatorname{GW}^{M}_{\{A_\gamma\},T}:
  H^*(M)^{\otimes k}\longrightarrow\mathbb Z
\]
为
\[
  \operatorname{GW}^{M}_{\{A_\gamma\},T}(a_1,\ldots,a_k)
  :=f\cdot\operatorname{ev}_{T,J,\{A_\gamma\}},
\]
这里 $f:U\to M^k$ 是与
$\pi_1^*a_1\smile\cdots\smile\pi_k^*a_k$ 对偶的\BookRecordIndex{T0337}伪循环。与
定理~\ref{thm:7-1-1} 的证明相同，这个数与 $J$ 无关。仅当各 $a_i$ 满足维数条件
\begin{equation}\label{eq:7-2-9}
  \sum_{i=1}^{k}\deg(a_i)
  =2n+2c_1(A)+2k-6-2e(T)
\end{equation}
时，它才可能非零。此外，定义
\[
  \operatorname{GW}^{M}_{A,T}(a_1,\ldots,a_k)
  :=\sum_{\substack{\{A_\gamma\}:\\
                     \sum_\gamma[A_\gamma]=A}}
    \operatorname{GW}^{M}_{\{A_\gamma\},T}(a_1,\ldots,a_k).
\]
从几何上说，不变量 $\operatorname{GW}^{M}_{A,T}(a_1,\ldots,a_k)$ 计数所有表示类 $A$、
以 $T$ 为模型的\BookRecordIndex{T0340}稳定映射 $(\mathbf u,\mathbf z)$，并且其第 $i$ 个\BookRecordIndex{T0066}标记点通过一个与
$a_i$ 对偶的循环 $Y_i$。

\begin{example}\label{ex:7-2-7}
令 $M=\mathbb{CP}^2$，令 $T$ 是有两个顶点 $0,1$ 的树，标号为
$\Lambda_0=\{1,2\}$、$\Lambda_1=\{3,4\}$，并令 $A_0=A_1=L$。给定四个互异点
$x_1,x_2,x_3,x_4\in\mathbb{CP}^2$，假设其中任意三点不共线。则存在唯一一对互异直线
$u_0,u_1:\mathbb{CP}^1\to\mathbb{CP}^2$，其中 $u_0$ 通过 $x_1,x_2$，而 $u_1$ 通过
$x_3,x_4$。由于两条直线的交点不等于任何 $x_i$，它们构成一个以 $T$ 为模型的\BookRecordIndex{T0340}稳定映射。
因此，用例~\ref{ex:7-1-14} 的记号，
\[
  \operatorname{GW}^{\mathbb{CP}^2}_{\{L,L\},T}
  (c^2,c^2,c^2,c^2)=1.
\]
注意，相应\BookRecordIndex{T0246}模空间的维数为 $16$，所以本例满足维数条件~\eqref{eq:7-2-9}。此外，$2L$ 的
其他分解 $\{0,2L\}$ 与 $\{2L,0\}$ 在这里都没有贡献，故
\[
  \operatorname{GW}^{\mathbb{CP}^2}_{2L,T}(c^2,c^2,c^2,c^2)=1.
\]
\end{example}

\begin{exercise}\label{ex:7-2-8}
选取上同调 $H^*(M)$ 的一组基 $\{e_0,\ldots,e_N\}$，并令 $g^{\nu\mu}$ 表示矩阵
\[
  g_{\nu\mu}:=\int_M e_\nu\smile e_\mu
\]
的逆矩阵。则上同调类
\[
  \sum_{\nu,\mu}e_\nu g^{\nu\mu}e_\mu
\]
与 $M\times M$ 的对角线庞加莱对偶。若 $T$ 与上一例相同，证明
\[
  \operatorname{GW}^{M}_{A,T}(a_1,\ldots,a_4)
  =\sum_{A_0+A_1=A}\sum_{\nu,\mu}
  \operatorname{GW}^{M}_{A_0,3}(a_1,a_2,e_\nu)
  g^{\nu\mu}
  \operatorname{GW}^{M}_{A_1,3}(e_\mu,a_3,a_4).
\]
\end{exercise}

\paragraph{\BookRecordIndex{T0150}固定标记点。}
下一节将定义\BookRecordIndex{T0143}格罗莫夫--威滕不变量 $\operatorname{GW}^{M,I}_{A,k}$，其中与指标集 $I$
相联系的\BookRecordIndex{T0066}标记点固定。这就是 Ruan--Tian \srcite{345} 所谓的\emph{\BookRecordIndex{T0172}混合不变量}。该定义
需要像第~\ref{sec:6-7} 节那样，用以底点 $z\in S^2$ 为参数的光滑\BookRecordIndex{T0088}殆复结构族来修正
Cauchy--Riemann 方程。作为准备，我们先用与 $z$ 无关的\BookRecordIndex{T0088}殆复结构讨论这类不变量的构造。
为此，必须对\BookRecordIndex{T0358}辛流形和指标集 $I$ 的基数施加相当强的限制。

假设 $(M,\omega)$ 单调，\BookRecordIndex{T0418}最小陈数为 $N$。固定满足条件~\eqref{eq:6-6-1} 的同调类
$A\in H_2(M;\mathbb Z)$、整数 $k\geq3$、子集 $I\subset\{1,\ldots,k\}$，使得
\[
  3\leq\#I\leq N+2,
\]
再固定 $S^2$ 上互异点所成的元组 $\mathbf w=\{w_i\}_{i\in I}$。回顾
习题~\ref{ex:6-6-5}：对 $J\in\mathcal J_{\mathrm{reg}}(M,\omega)$，\BookRecordIndex{T0283}求值映射
\[
  \operatorname{ev}_{\mathbf w,J}:
  \mathcal M^*_{0,k}(A,\mathbf w;J)\longrightarrow M^k
\]
是维数为 $2n+2c_1(A)+2(k-\#I)$ 的\BookRecordIndex{T0337}伪循环。因此，可以定义不变量
$\operatorname{GW}^{M,I}_{A,k}:H^*(M)^{\otimes k}\to\mathbb Z$：
\begin{equation}\label{eq:7-2-10}
  \operatorname{GW}^{M,I}_{A,k}(a_1,\ldots,a_k)
  :=f\cdot\operatorname{ev}_{\mathbf w,J}.
\end{equation}
这里 $f:U\to M^k$ 是与
$\pi_1^*a_1\smile\cdots\smile\pi_k^*a_k$ 对偶的\BookRecordIndex{T0337}伪循环，
$\pi_i:M^k\to M$ 是到第 $i$ 个因子的投影。与定理~\ref{thm:7-1-1} 的证明相同，这个数
与 $J$、$\mathbf w$ 都无关。仅当各 $a_i$ 满足维数条件
\begin{equation}\label{eq:7-2-11}
  \sum_{i=1}^{k}\deg(a_i)
  =2n+2c_1(A)+2k-2\#I
\end{equation}
时，它才可能非零。

从几何上说，不变量 $\operatorname{GW}^{M,I}_{A,k}(a_1,\ldots,a_k)$ 带适当符号计数
偶对
\[
  (u,\mathbf z)=\bigl(u,\{z_i\}_{i\notin I}\bigr),
\]
其中 $u:S^2\to M$ 是表示类 $A$ 的 $J$-全纯球面，而对
$i\in\{1,\ldots,k\}\setminus I$，各 $z_i$ 是
$S^2\setminus\{w_i\mid i\in I\}$ 中互异的点；此外，$i\in I$ 时
$u(w_i)\in X_i$，$i\notin I$ 时 $u(z_i)\in X_i$。所以，与 $I$ 中指标相联系的\BookRecordIndex{T0066}标记点
固定在 $S^2$ 上，其余\BookRecordIndex{T0066}标记点则可以自由变动。注意，上述定义对\BookRecordIndex{T0150}固定标记点的个数有很强的
限制 $\#I\leq N+2$。当 $N\geq2$ 时，它仍允许我们处理单调\BookRecordIndex{T0358}辛流形中重要的
$\#I=4$ 情形。

\begin{example}\label{ex:7-2-9}
为说明 $\operatorname{GW}$ 与 $\operatorname{GW}^{I}$ 的区别，考虑复射影平面
$M=\mathbb{CP}^2$ 中的二次曲线。注意，当 $A=2L$ 且所有 $a_i$ 都与点对偶时，
$I=\{1,\ldots,k\}$ 的 $\operatorname{GW}^{I}$ 所对应维数条件~\eqref{eq:7-2-11}
在 $k=4$ 时成立，而 $\operatorname{GW}$ 的维数条件~\eqref{eq:7-1-2} 在 $k=5$ 时
成立。经典地说，二次曲线是二次齐次多项式 $f$ 的零点集 $C_f$。五个一般位置点
（即任意三点不共线）唯一确定一条二次曲线这一熟知事实蕴含
\[
  \operatorname{GW}^{\mathbb{CP}^2}_{2L,5}(c^2,c^2,c^2,c^2,c^2)=1.
\]
（更多细节见习题~\ref{ex:7-1-16} 和定理~\ref{thm:7-4-1} 的证明。）于是，
$\mathbb{CP}^2$ 中通过四个固定一般位置点 $x_0,\ldots,x_3$ 的所有二次曲线构成实维数为
$2$ 的族。要从中选出唯一一条曲线，可以要求它通过 $\mathbb{CP}^2$ 中另一个点
（这给出不变量 $\operatorname{GW}$），也可以固定点 $x_0,\ldots,x_3$ 的原像
$w_0,\ldots,w_3$ 的\BookRecordIndex{T0185}交比。后一方法导出不变量
\[
  \operatorname{GW}^{\mathbb{CP}^2,\{1,\ldots,4\}}_{2L,4}
  (c^2,\ldots,c^2)
\]
的计算。

下面给出解析细节。给定 $w\in\mathbb C\setminus\{0,1\}$，存在唯一的二次全纯曲线
$u_w:\mathbb C\cup\{\infty\}\to\mathbb{CP}^2$，使得
\[
\begin{aligned}
  u_w(0)&=[1:0:0]=:x_0,
  &u_w(1)&=[0:1:0]=:x_1,\\
  u_w(\infty)&=[0:0:1]=:x_2,
  &u_w(w)&=[1:1:1]=:x_3,
\end{aligned}
\]
即
\[
  u_w(z)=[w(z-1):z(w-1):z(z-1)].
\]
此外，四元组 $(x_0,x_1,x_2,x_3)$ 是\BookRecordIndex{T0283}求值映射
\[
  \mathcal M^*(\mathbb{CP}^2,2L;J_0)
  \longrightarrow(\mathbb{CP}^2)^4,
  \qquad
  u\longmapsto\bigl(u(0),u(1),u(\infty),u(w)\bigr)
\]
的\BookRecordIndex{T0409}正则值。（这可以直接计算，也可以像命题~\ref{prop:7-4-5} 那样，考虑映射族
$g\circ u_w$，其中 $g\in\operatorname{PSL}(3,\mathbb C)$。）这说明
\[
  \operatorname{GW}^{\mathbb{CP}^2,\{1,\ldots,4\}}_{2L,4}
  (c^2,\ldots,c^2)=1.
\]

现在，对每个满足
$z_1\notin\{0,1\}$、$z_2\notin\{0,1\}$ 且 $z_1\ne z_2$ 的点
$x_4=[1:z_1:z_2]\in\mathbb{CP}^2$，存在唯一一对 $(w,z)$ 使 $u_w(z)=x_4$
\sourcefootnote{原书这里只要求 $z_1\ne1$、$z_2\ne0$，但这不足以保证前文所要求的
$w\notin\{0,1\}$ 以及 $z\notin\{0,1,\infty,w\}$。由显示公式直接化简，这些条件还要求
$z_1\ne0$、$z_2\ne1$、$z_1\ne z_2$；加上它们仍给出一个开稠密的 $x_4$ 集合，足以完成
一般点处的计数。}，即
\[
  w=\frac{z_1-z_2}{z_2(z_1-1)},
  \qquad
  z=\frac{z_1-z_2}{z_1-1}.
\]
而且，只要 $z\notin\{0,1,\infty,w\}$，映射
$(w,z)\mapsto u_w(z)$ 在 $(w,z)$ 处的微分便满射。因此，只要
$x_4=[1:z_1:z_2]$ 满足上述条件，点
$\mathbf x:=(x_0,x_1,x_2,x_3,x_4)$ 就是\BookRecordIndex{T0283}求值映射
\[
  \operatorname{ev}:\mathcal M^*_{0,5}(\mathbb{CP}^2,2L;J_0)
  \longrightarrow(\mathbb{CP}^2)^5
\]
的\BookRecordIndex{T0409}正则值，并且只有一个原像。故
\[
  \operatorname{GW}^{\mathbb{CP}^2}_{2L,5}(c^2,c^2,c^2,c^2,c^2)=1.
\]
\end{example}

\begin{remark}\label{rem:7-2-10}
上述论证表明
\[
  \operatorname{GW}^{\mathbb{CP}^2,\{1,\ldots,4\}}_{2L,4}
  (c^2,\ldots,c^2)=1.
\]
这个数等于例~\ref{ex:7-2-7} 中树 $T$ 的相应不变量
$\operatorname{GW}^{\mathbb{CP}^2}_{2L,T}(c^2,c^2,c^2,c^2)$；由
第~\ref{sec:6-7} 节最末关于完整\BookRecordIndex{T0145}格罗莫夫--威滕伪循环的讨论可知，该等式是粘接定理的
推论。
\end{remark}
